% Prova objetiva de Biblioteconomia
% Modelo do ÉguaTex (https://eguatex.github.io), em domínio público pela licença CC0 1.0:
% use, modifique e distribua sem pedir permissão e sem citar a origem.
%
% Compile com pdfLaTeX (o ÉguaTex faz isso sozinho). Arquivo único: main.tex.
% Classe exam: cada \question é uma questão, \choice é uma alternativa
% e \CorrectChoice marca a alternativa certa.

% ---------- versão do professor ou do aluno ----------
% Com \printanswers a alternativa certa sai destacada (gabarito).
% Troque por \noprintanswers para imprimir a prova do aluno.
\documentclass[11pt,a4paper,addpoints]{exam}
\printanswers

% ---------- idioma, fontes e página ----------
\usepackage[T1]{fontenc}
\usepackage[utf8]{inputenc}
\usepackage[brazilian]{babel}
\usepackage{tgpagella}
\usepackage{microtype}
\usepackage[a4paper,top=2.2cm,bottom=2.2cm,left=2.2cm,right=2.2cm,headheight=28pt]{geometry}
\usepackage{xcolor}
\usepackage{booktabs}
\usepackage[hidelinks]{hyperref}

\definecolor{vinho}{RGB}{122,30,48}

% ---------- instruções de preenchimento ----------
% Os textos em cinza explicam o que escrever em cada parte.
% Troque \instrucoestrue por \instrucoesfalse para esconder todos de uma vez.
\newif\ifinstrucoes
\instrucoestrue
\newcommand{\instrucao}[1]{\ifinstrucoes{\color{gray}\small\itshape #1\par}\fi}

% ---------- dados da prova: preencha aqui ----------
\newcommand{\universidade}{Universidade Federal Exemplo}
\newcommand{\curso}{Curso de Biblioteconomia}
\newcommand{\disciplina}{Fundamentos da Biblioteconomia}
\newcommand{\avaliacao}{Primeira avaliação}
\newcommand{\professor}{Nome do professor ou da professora}
\newcommand{\local}{Cidade}
\newcommand{\data}{Ano}

\hypersetup{pdftitle={\disciplina: \avaliacao}}

% ---------- termos da classe exam em português ----------
\pointpoints{ponto}{pontos}
\qformat{\textbf{Questão \thequestion}\hfill\textcolor{vinho}{\small(\thepoints)}}
\CorrectChoiceEmphasis{\color{vinho}\bfseries}
\renewcommand{\choicelabel}{\thechoice)}

% ---------- cabeçalho e rodapé ----------
\pagestyle{headandfoot}
\firstpageheader{}{}{}
\runningheader{\small\disciplina}{}{\small\avaliacao}
\runningheadrule
\footer{\small\universidade}{}{\small Página \thepage\ de \numpages}
\footrule

\begin{document}

% ======================================================================
% Identificação
% ======================================================================
\begin{center}
  {\large\bfseries\MakeUppercase{\universidade}\par}
  \curso\par
  \medskip
  {\Large\bfseries\color{vinho}\disciplina\par}
  \avaliacao\ \textbullet\ \local, \data\par
\end{center}

\vspace{1ex}
\noindent\begin{tabular}{@{}p{0.62\textwidth}p{0.32\textwidth}@{}}
  Nome: \hrulefill & Matrícula: \hrulefill \\[2ex]
  Professor(a): \professor & Nota: \hrulefill
\end{tabular}

\vspace{2ex}
\instrucao{Escreva as regras da prova em três a cinco itens curtos: duração, material permitido, como marcar a resposta e o valor das questões.}
\noindent\fbox{\parbox{\dimexpr\textwidth-2\fboxsep-2\fboxrule}{\small
  \textbf{Instruções.} A prova tem \numquestions\ questões objetivas e vale \numpoints\ pontos.
  Cada questão tem uma única alternativa correta. Marque a alternativa com caneta azul ou preta.
  Rasuras anulam a questão. Duração: 1 hora e 40 minutos.}}

\vspace{2ex}
\instrucao{Cada questão deve medir um conteúdo do plano de ensino. Escreva o enunciado completo e quatro alternativas plausíveis; troque \textbackslash choice por \textbackslash CorrectChoice na alternativa certa.}

% ======================================================================
% Questões
% ======================================================================
\begin{questions}

\question[1] Ranganathan publicou em 1931 as cinco leis da Biblioteconomia. Qual das frases abaixo \emph{não} é uma dessas leis?
\begin{choices}
  \choice Os livros são para usar.
  \choice A cada leitor o seu livro.
  \CorrectChoice Todo livro deve ser catalogado no dia em que chega.
  \choice Poupe o tempo do leitor.
\end{choices}

\question[1] A quinta lei de Ranganathan diz que:
\begin{choices}
  \choice a cada livro corresponde o seu leitor.
  \CorrectChoice a biblioteca é um organismo em crescimento.
  \choice o acervo deve ser descartado a cada dez anos.
  \choice a biblioteca deve priorizar obras de referência.
\end{choices}

\question[1] Na Classificação Decimal de Dewey (CDD), um livro de Física é classificado na classe principal:
\begin{choices}
  \choice 300, Ciências sociais.
  \CorrectChoice 500, Ciências naturais e Matemática.
  \choice 600, Tecnologia.
  \choice 700, Artes.
\end{choices}

\question[1] A Classificação Decimal Universal (CDU) foi desenvolvida a partir da CDD por:
\begin{choices}
  \choice Melvil Dewey e Charles Cutter.
  \choice S. R. Ranganathan.
  \CorrectChoice Paul Otlet e Henri La Fontaine.
  \choice Anthony Panizzi.
\end{choices}

\question[1] Desde 2007, o ISBN que identifica um livro tem:
\begin{choices}
  \choice 8 dígitos.
  \choice 10 dígitos.
  \CorrectChoice 13 dígitos.
  \choice 16 dígitos.
\end{choices}

\question[1] O ISSN é o identificador padronizado de:
\begin{choices}
  \CorrectChoice publicações seriadas, como periódicos e revistas.
  \choice capítulos de livros.
  \choice teses e dissertações.
  \choice registros de autoridade de nomes.
\end{choices}

\question[1] A norma da ABNT que estabelece como apresentar citações em documentos é a:
\begin{choices}
  \choice NBR 6023.
  \choice NBR 14724.
  \CorrectChoice NBR 10520.
  \choice NBR 6028.
\end{choices}

\question[1] No formato MARC 21 bibliográfico, o campo que registra o título principal da obra é o:
\begin{choices}
  \choice 020.
  \choice 100.
  \CorrectChoice 245.
  \choice 650.
\end{choices}

\question[1] O código de catalogação que substituiu o AACR2 e se baseia nos modelos conceituais FRBR é o:
\begin{choices}
  \choice ISBD.
  \CorrectChoice RDA (Resource Description and Access).
  \choice Dublin Core.
  \choice CCAA.
\end{choices}

\question[1] Na Classificação de Dois Pontos, de Ranganathan, a fórmula PMEST reúne as categorias fundamentais:
\begin{choices}
  \choice Público, Meio, Edição, Suporte e Título.
  \CorrectChoice Personalidade, Matéria, Energia, Espaço e Tempo.
  \choice Proveniência, Método, Estrutura, Série e Tema.
  \choice Período, Material, Escola, Sistema e Tradição.
\end{choices}

\end{questions}

% ======================================================================
% Quadro de respostas e gabarito
% ======================================================================
\vspace{3ex}
\instrucao{O quadro abaixo facilita a correção. Na versão do professor ele traz o gabarito; na do aluno, fica em branco.}
% \resp{x} mostra a letra só na versão do professor.
\newcommand{\resp}[1]{\ifprintanswers\textcolor{vinho}{\textbf{#1}}\fi}
\begin{center}
  \renewcommand{\arraystretch}{1.4}
  \begin{tabular}{l*{10}{c}}
    \toprule
    Questão & 1 & 2 & 3 & 4 & 5 & 6 & 7 & 8 & 9 & 10 \\
    \midrule
    Resposta & \resp{c} & \resp{b} & \resp{b} & \resp{c} & \resp{c} & \resp{a} & \resp{c} & \resp{c} & \resp{b} & \resp{b} \\
    \bottomrule
  \end{tabular}
\end{center}

\end{document}
